\documentclass[tikz, border=6pt]{standalone}
\usepackage[fontset=fandol]{ctex}
\usetikzlibrary{positioning, arrows.meta, calc, fit, backgrounds}

\definecolor{cMain}{RGB}{62, 96, 138}
\definecolor{cWarm}{RGB}{176, 118, 62}
\definecolor{cGray}{RGB}{120, 120, 120}

% #1 节点名  #2 颜色  #3 顶部中心坐标  #4 宽度  #5 标题  #6 说明
\newcommand{\card}[6]{%
    \node[draw=#2!80!black, fill=#2!14, line width=0.6pt, minimum width=#4,
          minimum height=0.62cm, inner sep=2pt, font=\heiti\zihao{-5}, anchor=north] (#1-t) at (#3) {#5};
    \node[draw=#2!80!black, fill=white, line width=0.6pt, minimum width=#4,
          minimum height=1.55cm, inner sep=3pt, align=center, anchor=north,
          font=\zihao{-5}] (#1-b) at (#1-t.south) {#6};
    \node[fit=(#1-t)(#1-b), inner sep=0pt] (#1) {};
}

\begin{document}
\begin{tikzpicture}[
    arr/.style={-{Stealth[length=2.2mm, width=1.6mm]}, line width=0.6pt, draw=black!75},
]

\def\wS{2.9cm}
\def\wE{3.6cm}
\def\yA{0}
\def\yB{-3.0}

% ---------------- 原始数据 ----------------
\card{raw}{cGray}{0, \yA}{\wE}{原始数据}{内部语料 / 公开语料\\PDF / Word\\表格 / 工单}

% ---------------- 预处理步骤 1--6（两行回转） ----------------
\card{s1}{cMain}{4.25, \yA}{\wS}{1. 文本抽取}{差异化文档解析\\双栏与表格还原\\统一 UTF-8 文本}
\card{s2}{cMain}{7.85, \yA}{\wS}{2. 格式规范}{统一条款层级\\标准化计量单位\\规范专业符号}
\card{s3}{cMain}{11.45, \yA}{\wS}{3. 噪声清洗}{剔除页眉与页码\\清理 OCR 乱码\\过滤重复声明}
\card{s4}{cMain}{11.45, \yB}{\wS}{4. 分段切片}{4096 token 窗口\\自然句边界截断\\相邻逻辑块拼接}
\card{s5}{cMain}{7.85, \yB}{\wS}{5. 去重筛选}{一级 SimHash\\二级 MinHash\\三级 BGE-M3}
\card{s6}{cMain}{4.25, \yB}{\wS}{6. 质量评估}{自动指标初筛\\LLM 辅助判断\\人工抽样校验}

% ---------------- 标签化 ----------------
\card{lab}{cWarm}{0, \yB}{\wE}{7. 七类数据类型标签化}{用于 CPT 知识溯源、\\SFT 场景原生指令生成\\与评测文档分层抽取}

% ---------------- 连线 ----------------
\draw[arr] (raw) -- (s1);
\draw[arr] (s1) -- (s2);
\draw[arr] (s2) -- (s3);
\draw[arr] (s3) -- (s4);
\draw[arr] (s4) -- (s5);
\draw[arr] (s5) -- (s6);
\draw[arr] (s6) -- (lab);

% ---------------- 阶段背景 ----------------
\begin{scope}[on background layer]
    \node[fit=(s1)(s3)(s4)(s6), draw=cMain!45, fill=cMain!4, line width=0.5pt,
          inner xsep=0.22cm, inner ysep=0.22cm] (bgA) {};
\end{scope}
\node[anchor=south west, inner sep=2pt, font=\zihao{-5}, text=cMain!90!black]
    at (bgA.north west) {阶段 A：系统化预处理（步骤 1--6）};
\node[anchor=south west, inner sep=2pt, font=\zihao{-5}, text=cWarm!90!black]
    at ([yshift=0.22cm]lab.north west) {阶段 B：标签化（步骤 7）};

\end{tikzpicture}
\end{document}
